\subsection{限制与余限制}

设 G 与 \(G'\) 是群，F 是交换群，而 \(\tau\) 是从 G 到 F 的自同构群的一个同态。设 \(u: G' \to G\) 是一个群同态。我们记 \(\tau' = \tau \circ u\)。

若 \(\psi : \mathrm{G} \times \mathrm{G} \to \mathrm{F}\) 是 \(\mathrm{G}\) 的取值于 \(\mathrm{F}\) 的一个 2-闭上链，则映射 \(\psi \circ (u \times u)\)（从 \(\mathrm{G}' \times \mathrm{G}'\) 到 \(\mathrm{F}\)，由 \((g_1, g_2) \mapsto \psi(u(g_1), u(g_2))\) 给出）是 \(\mathrm{G}'\) 的取值于 \(\mathrm{F}\) 的一个 2-闭上链，且映射 \(\psi \mapsto \psi \circ (u \times u)\)（从 \(\mathrm{Z}^2(\mathrm{G}, \mathrm{F})\) 到 \(\mathrm{Z}^2(\mathrm{G}', \mathrm{F})\)）诱导一个群同态 \(u^*: \mathrm{H}^2(\mathrm{G}, \mathrm{F}) \to \mathrm{H}^2(\mathrm{G}', \mathrm{F})\)。对任何 \(\lambda \in \mathrm{H}^2(\mathrm{G}, \mathrm{F})\)，元素 \(u^*(\lambda)\) 称为 \(\lambda\) 被 \(u\) 的逆像。当 \(\mathrm{G}'\) 是 \(\mathrm{G}\) 的子群且 \(u: \mathrm{G}' \to \mathrm{G}\) 是典范单射时，同态 \(u^*\) 称为从 \(\mathrm{G}\) 到 \(\mathrm{G}'\) 的限制同态；我们把它记作 \(\operatorname{Res}_{\mathrm{G}'}^{\mathrm{G}}\)。当 \(\mathrm{G}\) 是 \(\mathrm{G}'\) 的商且 \(u: \mathrm{G}' \to \mathrm{G}\) 是典范满射时，同态 \(u^*\) 称为从 \(\mathrm{G}\) 到 \(\mathrm{G}'\) 的膨胀同态。

命题 7. —— 用上面的记号，下图交换：

\[
\begin{array}{c} \operatorname{Ex} _ {\tau} (\mathrm{G}, \mathrm{F}) \xrightarrow {u ^ {*}} \operatorname{Ex} _ {\tau} (\mathrm{G} ^ {\prime}, \mathrm{F}) \\ \Biggl \downarrow \Theta_ {\tau} \qquad \qquad \qquad \Biggl \downarrow \Theta_ {\tau^ {\prime}} \\ \operatorname{H} ^ {2} (\mathrm{G}, \mathrm{F}) \xrightarrow {u ^ {*}} \operatorname{H} ^ {2} (\mathrm{G} ^ {\prime}, \mathrm{F})  . \end{array}
\]

设 H 是 G 的一个有限指数子群。设 s 是从 G 到 \(H\backslash G\) 的典范满射的一个截面。我们用 \((g,x)\mapsto x\cdot g\) 表示由 G 借乘法在其自身上的右作用所诱导的 G 在 \(H\backslash G\) 上的右作用。注意对每个 \(x\in H\backslash G\) 与 \(g\in G\)，元素 \(s(x)gs(x\cdot g)^{-1}\) 属于 H。于是对任何映射 \(c: \mathrm{H} \times \mathrm{H} \to \mathrm{F}\)，我们都可以用关系式定义一个映射 \(\widetilde{c}_s: \mathrm{G} \times \mathrm{G} \to \mathrm{F}\)

\[
\widetilde {c} _ {s} (g _ {1}, g _ {2}) = \prod_ {x \in \mathrm{H} \backslash \mathrm{G}} ^ {s (x) ^ {- 1}} c \bigl (s (x) g _ {1} s (x \cdot g _ {1}) ^ {- 1}, s (x \cdot g _ {1}) g _ {2} s (x \cdot g _ {1} g _ {2}) ^ {- 1} \bigr)
\]

对 \(g_{1}, g_{2} \in G\)。映射 \(c \mapsto \widetilde{c}_{s}\) 是从 \(F^{H \times H}\) 到 \(F^{G \times G}\) 的一个群同态。

引理 4. —— 若 \(c: \mathrm{H} \times \mathrm{H} \to \mathrm{F}\) 是 H 的取值于 F 的一个 2-闭上链，则 \(\widetilde{c}_s\) 是 G 的取值于 F 的一个 2-闭上链。

设 \(g_1, g_2\) 与 \(g_3\) 是 G 的元素。对每个 \(i \in \{1, 2, 3\}\)，我们用关系式定义一个映射 \(h_i: \mathrm{H} \backslash \mathrm{G} \to \mathrm{H}\)

\[
h _ {i} (x) = s \left(x \cdot g _ {1} \dots g _ {i - 1}\right) g _ {i} s \left(x \cdot g _ {1} \dots g _ {i}\right) ^ {- 1}
\]

对 \(x\in \mathrm{H}\backslash \mathrm{G}\)。注意

\[
\begin{array}{l} {h _ {1} (x) h _ {2} (x) = s (x) g _ {1} g _ {2} s (x \cdot g _ {1} g _ {2}) ^ {- 1}} \\ {h _ {2} (x) h _ {3} (x) = s (x \cdot g _ {1}) g _ {2} g _ {3} s (x \cdot g _ {1} g _ {2} g _ {3}) ^ {- 1}} \end{array}
\]

对 \(x \in \mathrm{H} \backslash \mathrm{G}\)。于是我们有关系

\[
\begin{array}{r l} & ^ {g _ {1}} \widetilde {c} _ {s} (g _ {2}, g _ {3}) \widetilde {c} _ {s} (g _ {1} g _ {2}, g _ {3}) ^ {- 1} \widetilde {c} _ {s} (g _ {1}, g _ {2} g _ {3}) \widetilde {c} _ {s} (g _ {1}, g _ {2}) ^ {- 1} \\ & \quad = \prod_ {x \in \mathrm{H} \backslash \mathrm{G}} ^ {g _ {1} s (x) ^ {- 1}} c \big (s (x) g _ {2} s (x \cdot g _ {2}) ^ {- 1}, s (x \cdot g _ {2}) g _ {3} s (x \cdot g _ {2} g _ {3}) ^ {- 1} \big) \\ & \qquad \times \prod_ {x \in \mathrm{H} \backslash \mathrm{G}} ^ {s (x) ^ {- 1}} \Big (c (h _ {1} (x) h _ {2} (x), h _ {3} (x)) ^ {- 1} \\ & \qquad \qquad \qquad \qquad \times c (h _ {1} (x), h _ {2} (x) h _ {3} (x))   c (h _ {1} (x), h _ {2} (x)) ^ {- 1} \Big) \\ & = \prod_ {x \in \mathrm{H} \backslash \mathrm{G}} ^ {s (x \cdot g _ {1} ^ {- 1}) ^ {- 1} h _ {1} (x \cdot g _ {1} ^ {- 1})} c (h _ {2} (x \cdot g _ {1} ^ {- 1}), h _ {3} (x \cdot g _ {1} ^ {- 1})) \\ & \qquad \times \prod_ {x \in \mathrm{H} \backslash \mathrm{G}} ^ {s (x) ^ {- 1} h _ {1} (x)} c (h _ {2} (x), h _ {3} (x)) ^ {- 1} \\ & = 1, \end{array}
\]

其中第二个等式由 c 是 2-闭上链这一事实得出。

引理 5. —— 若 \(c\) 是一个 2-上边缘，则 \(\widetilde{c}_s\) 也是。

设 \(t: \mathrm{H} \to \mathrm{F}\) 是使得 \(c = \partial t\) 的一个映射。设 \(\widetilde{t}_s: \mathrm{G} \to \mathrm{F}\) 是由下式定义的映射

\[
\widetilde {t} _ {s} (g) = \prod_ {x \in \mathrm{H} \backslash \mathrm{G}} ^ {s (x) ^ {- 1}} t (s (x) g s (x \cdot g) ^ {- 1})
\]

对 \(g \in G\)。只需证明 \(\widetilde{c}_{s} = \partial\widetilde{t}_{s}\)，这由关系

\[
\begin{array}{l} \widetilde {c} _ {s} (g _ {1}, g _ {2}) = \prod_ {x \in \mathrm{H} \backslash \mathrm{G}} ^ {s (x) ^ {- 1} s (x) g _ {1} s (x \cdot g _ {1}) ^ {- 1}} t (s (x \cdot g _ {1}) g _ {2} s (x \cdot g _ {1} g _ {2}) ^ {- 1}) \\ \qquad \times \prod_ {x \in \mathrm{H} \backslash \mathrm{G}} ^ {s (x) ^ {- 1}} \big (t (s (x) g _ {1} g _ {2} s (x \cdot g _ {1} g _ {2}) ^ {- 1}) ^ {- 1} t (s (x) g _ {1} s (x \cdot g _ {1}) ^ {- 1}) \big) \\ = \partial \widetilde {t} _ {s} (g _ {1}, g _ {2}) \end{array}
\]

对 \(g_{1}, g_{2} \in G\) 得出。

引理 6. —— 设 c 是 H 的取值于 F 的一个 2-闭上链。\(\widetilde{c}_{s}\) 在群 \(\mathrm{H}^{2}(\mathrm{G},\mathrm{F})\) 中的像不依赖于截面 s 的选取。

设 \(s'\) 是从 G 到 \(H\backslash G\) 的典范满射的一个截面。设 \(h: \mathrm{H} \backslash \mathrm{G} \to \mathrm{H}\) 是由关系式

\[
s ^ {\prime} (x) = h (x) s (x)
\]

对 \(x \in H \backslash G\) 刻画的映射。于是 2-闭上链 \(\widetilde{c}_{s'}\) 由关系式给出

\[
\begin{array}{l} \widetilde {c} _ {s ^ {\prime}} (g _ {1}, g _ {2}) \\ = \prod_ {x \in \mathrm{H} \backslash \mathrm{G}} ^ {s (x) ^ {- 1} h (x) ^ {- 1}} c \big (h (x) s (x) g _ {1} s ^ {\prime} (x \cdot g _ {1}) ^ {- 1}, s ^ {\prime} (x \cdot g _ {1}) g _ {2} s ^ {\prime} (x \cdot g _ {1} g _ {2}) ^ {- 1} \big) \\ = \prod_ {x \in \mathrm{H} \backslash \mathrm{G}} ^ {s (x) ^ {- 1}} c \big (s (x) g _ {1} s (x \cdot g _ {1}) ^ {- 1} h (x \cdot g _ {1}) ^ {- 1}, h (x \cdot g _ {1}) s (x \cdot g _ {1}) g _ {2} s ^ {\prime} (x \cdot g _ {1} g _ {2}) ^ {- 1} \big) \\ \qquad \times \prod_ {x \in \mathrm{H} \backslash \mathrm{G}} ^ {s (x) ^ {- 1}} \Big (c \big (h (x) ^ {- 1}, s ^ {\prime} (x) g _ {1} g _ {2} s ^ {\prime} (x \cdot g _ {1} g _ {2}) ^ {- 1} \big) ^ {- 1} \\ \qquad \qquad \qquad \qquad \qquad \qquad \qquad \qquad \qquad \qquad \qquad \qquad \qquad \qquad \qquad \qquad \qquad \qquad \qquad \qquad \qquad \qquad \qquad \qquad \qquad \qquad \qquad \qquad \qquad \qquad \qquad \qquad \qquad \qquad \\ = \prod_ {x \in \mathrm{H} \backslash \mathrm{G}} ^ {g _ {1} s (x \cdot g _ {1}) ^ {- 1}} c \big (h (x \cdot g _ {1}) ^ {- 1}, h (x \cdot g _ {1}) s (x \cdot g _ {1}) g _ {2} s ^ {\prime} (x \cdot g _ {1} g _ {2}) ^ {- 1} \big) \\ \qquad \times \prod_ {x \in \mathrm{H} \backslash \mathrm{G}} ^ {s (x) ^ - 1} c \big (s (x) g _ {1} s (x \cdot g _ {1}) ^ {- 1}, s (x \cdot g _ {1}) g _ {2} s ^ {\prime} (x \cdot g _ {1} g _ {2}) ^ {- 1} \big) \\ \qquad \times \prod_ {x \in \mathrm{H} \backslash \mathrm{G}} ^ {s (x) ^ {- 1}} c \big (s (x) g _ {1} s (x \cdot g _ {1}) ^ {- 1}, h (x \cdot g _ {1}) ^ {- 1} \big) ^ {- 1} \\ \qquad \times \prod_ {x \in \mathrm{H} \backslash \mathrm{G}} ^ {s (x) ^ {- 1}} \Big (c \big (h (x) ^ {- 1}, s ^ {\prime} (x) g _ {1} g _ {2} s ^ {\prime} (x \cdot g _ {1} g _ {2}) ^ {- 1} \big) ^ {{- 1}} \\ \qquad \qquad \qquad \qquad \qquad \qquad \qquad \qquad \qquad \qquad \qquad \qquad \times c (h (x) ^ {- 1}, s ^ {\prime} (x) g _ {1} s ^ {\prime} (x \cdot g _ {1}) ^ {- 1})   {\Big)} \end{array}
\]

\[
\begin{array}{l} = \prod_ {x \in \mathrm{H} \backslash \mathrm{G}} ^ {s (x) ^ {- 1}} c \big (s (x) g _ {1} s (x \cdot g _ {1}) ^ {- 1}, s (x \cdot g _ {1}) g _ {2} s ^ {\prime} (x \cdot g _ {1} g _ {2}) ^ {- 1} \big) \\ \qquad \times \prod_ {x \in \mathrm{H} \backslash \mathrm{G}} ^ {s (x) ^ {- 1}} c \big (s (x) g _ {1} s (x \cdot g _ {1}) ^ {- 1}, h (x \cdot g _ {1}) ^ {- 1} \big) ^ {- 1} \\ \qquad \times \prod_ {x \in \mathrm{H} \backslash \mathrm{G}} ^ {g _ {1} s (x) ^ {- 1}} c \big (h (x) ^ {- 1}, h (x) s (x) g _ {2} s ^ {\prime} (x \cdot g _ {2}) ^ {- 1} \big) \\ \qquad \times \prod_ {x \in \mathrm{H} \backslash \mathrm{G}} ^ {s (x) ^ {- 1}} \Big (c \big (h (x) ^ {- 1}, s ^ {\prime} (x) g _ {1} g _ {2} s ^ {\prime} (x \cdot g _ {1} g _ {2}) ^ {- 1} \big) ^ {- 1} \\ \qquad \qquad \qquad \qquad \qquad \qquad \qquad \qquad \qquad \qquad \qquad \qquad \qquad \qquad \qquad \qquad \qquad \qquad \times c \big (h (x) ^ {- 1}, s ^ {\prime} (x) g _ {1} s ^ {\prime} (x \cdot g _ {1}) ^ {- 1} \big) \Big) \end{array}
\]

对所有 \(g_1, g_2 \in \mathbf{G}\)。第一个等式来自把闭上链关系（VIII, 第 295 页, 公式 (7)）应用于元素

\[
h (x) ^ {- 1}, \quad h (x) s (x) g _ {1} s ^ {\prime} (x \cdot g _ {1}) ^ {- 1}, \quad \mathrm{and} \quad s ^ {\prime} (x \cdot g _ {1}) g _ {2} s ^ {\prime} (x \cdot g _ {1} g _ {2}) ^ {- 1};
\]

第二个等式是通过对元素

\[
s (x) g _ {1} s \left(x \cdot g _ {1}\right) ^ {- 1}, \quad h \left(x \cdot g _ {1}\right) ^ {- 1}, \quad \text { and } \quad h \left(x \cdot g _ {1}\right) s \left(x \cdot g _ {1}\right) g _ {2} s ^ {\prime} \left(x \cdot g _ {1} g _ {2}\right) ^ {- 1};
\]

而最后一个等式只是用到了映射 \(x \mapsto x \cdot g_{1}\) 是 \(H\backslash G\) 的一个置换这一事实。

所得表达式的最后两行对应于一个 2-上边缘。我们发现 \(\widetilde{c}_{s'}\) 在 \(\mathrm{H}^2 (\mathrm{G},\mathrm{F})\) 中与那个在 \((g_{1},g_{2})\in \mathrm{G}^{2}\) 处的值由下述表达式给出的闭上链有相同的类

\[
\begin{array}{l} \prod_ {x \in \mathrm{H} \backslash \mathrm{G}} ^ {s (x) ^ {- 1}} c \big (s (x) g _ {1} s (x \cdot g _ {1}) ^ {- 1}, s (x \cdot g _ {1}) g _ {2} s (x \cdot g _ {1} g _ {2}) ^ {- 1} h (x \cdot g _ {1} g _ {2}) ^ {- 1} \big) \\ \qquad \qquad \qquad \times \prod_ {x \in \mathrm{H} \backslash \mathrm{G}} ^ {s (x) ^ {- 1}} c \big (s (x) g _ {1} s (x \cdot g _ {1}) ^ {- 1}, h (x \cdot g _ {1}) ^ {- 1} \big) ^ {- 1} \\ = \prod_ {x \in \mathrm{H} \backslash \mathrm{G}} ^ {g _ {1} s (x \cdot g _ {1}) ^ {- 1}} c \big (s (x \cdot g _ {1}) g _ {2} s (x \cdot g _ {1} g _ {2}) ^ {- 1}, h (x \cdot g _ {1} g _ {2}) ^ {- 1} \big) ^ {- 1} \\ \qquad \qquad \times \prod_ {x \in \mathrm{H} \backslash \mathrm{G}} ^ {s (x) ^ {- 1}} c \big (s (x) g _ {1} s (x \cdot g _ {1}) ^ {- 1}, s (x \cdot g _ {1}) g _ {2} s (x \cdot g _ {1} g _ {2}) ^ {- 1} \big) \\ \qquad \qquad \times \prod_ {x \in \mathrm{H} \backslash \mathrm{G}} ^ {s (x) ^ {- 1}} \Big (c \big (s (x) g _ {1} g _ {2} s (x \cdot g _ {1} g _ {2}) ^ {- 1}, h (x \cdot g _ {1} g _ {2}) ^ {- 1} \big) \\ \qquad \qquad \qquad \times c \big (s (x) g _ {1} s (x \cdot g _ {1}) ^ {- 1}, h (x \cdot g _ {1}) ^ {- 1} \big) ^ {- 1} \Big) \end{array}
\]

\[
\begin{array}{l} = \prod_ {x \in \mathrm{H} \backslash \mathrm{G}} ^ {s (x) ^ {- 1}} c \big (s (x) g _ {1} s (x \cdot g _ {1}) ^ {- 1}, s (x \cdot g _ {1}) g _ {2} s (x \cdot g _ {1} g _ {2}) ^ {- 1} \big) \\ \quad \times \prod_ {x \in \mathrm{H} \backslash \mathrm{G}} ^ {g _ {1} s (x) ^ {- 1}} c \big (s (x) g _ {2} s (x \cdot g _ {2}) ^ {- 1}, h (x \cdot g _ {2}) ^ {- 1} \big) ^ {- 1} \\ \quad \times \prod_ {x \in \mathrm{H} \backslash \mathrm{G}} ^ {s (x) ^ {- 1}} \Big (c \big (s (x) g _ {1} g _ {2} s (x \cdot g _ {1} g _ {2}) ^ {- 1}, h (x \cdot g _ {1} g _ {2}) ^ {- 1} \big) \\ \qquad \qquad \qquad \qquad \qquad \qquad \qquad \qquad \qquad \qquad \qquad \qquad \qquad \times c \big (s (x) g _ {1} s (x \cdot g _ {1}) ^ {- 1}, h (x \cdot g _ {1}) ^ {- 1} \big) ^ {- 1} \Big), \end{array}
\]

其中第一个等式由对元素应用闭上链关系得出

\[
s (x) g _ {1} s (x \cdot g _ {1}) ^ {- 1}, \quad s (x \cdot g _ {1}) g _ {2} s (x \cdot g _ {1} g _ {2}) ^ {- 1}, \quad \text { and } \quad h (x \cdot g _ {1} g _ {2}) ^ {- 1}.
\]

所得表达式的最后两行对应于一个 2-上边缘。我们发现 \(\widetilde{c}_{s'}\) 的类与 \(\widetilde{c}_s\) 的类重合。

我们已经构造了一个从群 \(\mathrm{H}^{2}(\mathrm{H},\mathrm{F})\) 到群 \(\mathrm{H}^{2}(\mathrm{G},\mathrm{F})\) 的同态，我们称它为从 H 到 G 的余限制同态，并记作 \(Cor_{H}^{G}\)。

命题 8. —— 设 H 是 G 的一个有限指数子群。自同态 \(Cor_{H}^{G} \circ Res_{H}^{G}\)（群 \(H^{2}(G, F)\) 的）与用指数 \((G : H)\) 作乘法相重合。

设 \(\alpha\) 是 \(\mathrm{H}^2 (\mathrm{G},\mathrm{F})\) 的一个元素，并设 \(c\) 是 \(\mathrm{Z}^2 (\mathrm{G},\mathrm{F})\) 的一个元素，它代表 \(\alpha\)。元素 \(\mathrm{Cor}_{\mathrm{H}}^{\mathrm{G}}\circ \mathrm{Res}_{\mathrm{H}}^{\mathrm{G}}(\alpha)\) 是那个在 \((g_{1},g_{2})\in \mathrm{G}^{2}\) 处的值由下述表达式给出的闭上链的类

\[
\begin{array}{l} \prod_ {x \in \mathrm{H} \backslash \mathrm{G}} ^ {s (x) ^ {- 1}} c (s (x) g _ {1} s (x \cdot g _ {1}) ^ {- 1}, s (x \cdot g _ {1}) g _ {2} s (x \cdot g _ {1} g _ {2}) ^ {- 1}) \\ = \prod_ {x \in \mathrm{H} \backslash \mathrm{G}} c (g _ {1} s (x \cdot g _ {1}) ^ {- 1}, s (x \cdot g _ {1}) g _ {2} s (x \cdot g _ {1} g _ {2}) ^ {- 1}) \\ \quad \times \prod_ {x \in \mathrm{H} \backslash \mathrm{G}} \big (c (s (x) ^ {- 1}, s (x) g _ {1} g _ {2} s (x \cdot g _ {1} g _ {2}) ^ {- 1}) ^ {- 1} c (s (x) ^ {- 1}, s (x) g _ {1} s (x \cdot g _ {1}) ^ {- 1}) \big) \\ = \prod_ {x \in \mathrm{H} \backslash \mathrm{G}} ^ {g _ {1}} c (s (x \cdot g _ {1}) ^ {- 1}, s (x \cdot g _ {1}) g _ {2} s (x \cdot g _ {1} g _ {2}) ^ {- 1}) \\ \quad \times \prod_ {x \in \mathrm{H} \backslash \mathrm{G}} c (g _ {1}, g _ {2} s (x \cdot g _ {1} g _ {2}) ^ {- 1}) c (g _ {1}, s (x \cdot g _ {1}) ^ {- 1}) ^ {- 1} \\ \quad \times \prod_ {x \in \mathrm{H} \backslash \mathrm{G}} \big (c (s (x) ^ {- 1}, s (x) g _ {1} g _ {2} s (x \cdot g _ {1} g _ {2}) ^ {- 1}) ^ {- 1} c (s (x) ^ {- 1}, s (x) g _ {1} s (x \cdot g _ {1}) ^ {- 1}) \big) \end{array}
\]

\[
\begin{array}{l} = \prod_ {x \in \mathrm{H} \backslash \mathrm{G}} c (g _ {1}, g _ {2} s (x \cdot g _ {1} g _ {2}) ^ {- 1}) c (g _ {1}, s (x \cdot g _ {1}) ^ {- 1}) ^ {- 1} \\ \times \prod_ {x \in \mathrm{H} \backslash \mathrm{G}} ^ {g _ {1}} c (s (x) ^ {- 1}, s (x) g _ {2} s (x g _ {2}) ^ {- 1}) \\ \times \prod_ {x \in \mathrm{H} \backslash \mathrm{G}} \left(c (s (x) ^ {- 1}, s (x) g _ {1} g _ {2} s (x \cdot g _ {1} g _ {2}) ^ {- 1}) ^ {- 1} c (s (x) ^ {- 1}, s (x) g _ {1} s (x \cdot g _ {1}) ^ {- 1})\right). \end{array}
\]

第一个等式来自把闭上链关系（VIII, 第 295 页, 公式 (7)）应用于元素

\[
s (x) ^ {- 1}, \quad s (x) g _ {1} s (x \cdot g _ {1}) ^ {- 1}, \quad \mathrm{and} \quad s (x \cdot g _ {1}) g _ {2} s (x \cdot g _ {1} g _ {2}) ^ {- 1};
\]

第二个等式是通过对元素

\[
g _ {1}, \quad s (x \cdot g _ {1}) ^ {- 1}, \quad \mathrm{and} \quad s (x \cdot g _ {1}) g _ {2} s (x \cdot g _ {1} g _ {2}) ^ {- 1}.
\]

通过去掉一个上边缘，我们得到 \(\mathrm{Cor}_{\mathrm{H}}^{\mathrm{G}} \circ \mathrm{Res}_{\mathrm{H}}^{\mathrm{G}}(\alpha)\) 是那个在 \((g_{1}, g_{2}) \in \mathrm{G}^{2}\) 处的值由下述表达式给出的闭上链的类

\[
\begin{array}{l} \prod_ {x \in \mathrm{H} \backslash \mathrm{G}} c (g _ {1}, g _ {2} s (x \cdot g _ {1} g _ {2}) ^ {- 1}) c (g _ {1}, s (x \cdot g _ {1}) ^ {- 1}) ^ {- 1} \\ = \prod_ {x \in \mathrm{H} \backslash \mathrm{G}} ^ {g _ {1}} c (g _ {2}, s (x \cdot g _ {1} g _ {2}) ^ {- 1}) ^ {- 1} c (g _ {1} g _ {2}, s (x \cdot g _ {1} g _ {2}) ^ {- 1}) c (g _ {1}, g _ {2}) \\ \times \prod_ {x \in \mathrm{H} \backslash \mathrm{G}} c (g _ {1}, s (x \cdot g _ {1}) ^ {- 1}) ^ {- 1} \\ = \prod_ {x \in \mathrm{H} \backslash \mathrm{G}} ^ {g _ {1}} c (g _ {2}, s (x \cdot g _ {2}) ^ {- 1}) ^ {- 1} c (g _ {1} g _ {2}, s (x \cdot g _ {1} g _ {2}) ^ {- 1}) c (g _ {1}, s (x \cdot g _ {1}) ^ {- 1}) ^ {- 1} \\ \times c (g _ {1}, g _ {2}) ^ {\text {(G:H)}}, \end{array}
\]

这就证明了结论。

